\documentclass[12pt]{article}
\usepackage[margin=1in]{geometry} 
\usepackage{amsmath,amsthm,amssymb,scrextend}
\setlength{\parindent}{0pt}
\begin{document}


\textbf{(i)} Let $u' = x^n \;\mathrm{d}x$, $v = \arctan x$

$\implies u=\dfrac{1}{n+1}\;x^{n+1}$, $v = \dfrac{1}{1+x^2}\,\mathrm{d}x$

$$\displaystyle\int uv'\mathrm{d}x=uv-\displaystyle\int u'v\mathrm{d}x$$

\begin{align*}
I_n &= \left[\frac{x^{n+1}}{n+1}\arctan x\right]_0^1 - \int_0^1 \frac{x^{n+1}}{n+1}\cdot\frac{1}{1+x^2}\,\mathrm{d}x\\
&= \frac{1}{n+1}\;\frac{\pi}{4} - \frac{1}{n+1}\int_0^1 \frac{x^{n+1}}{1+x^2}\,\mathrm{d}x\\
(n+1)\;I_n &= \frac{\pi}{4} - \int_0^1 \frac{x^{n+1}}{1+x^2}\,\mathrm{d}x
\end{align*}
When $n=0$, $\qquad I_0 = \dfrac{\pi}{4} - \displaystyle\int_0^1 \frac{x}{1+x^2}\,\mathrm{d}x = \frac{\pi}{4} - \frac{1}{2}\ln 2$



\bigskip\hrule\bigskip
\textbf{(ii)} From part (i):
$$(n+3)I_{n+2} = \frac{\pi}{4} - \int_0^1 \frac{x^{n+3}}{1+x^2}\,\mathrm{d}x, \qquad (n+1)I_n = \frac{\pi}{4} - \int_0^1 \frac{x^{n+1}}{1+x^2}\,\mathrm{d}x$$

\begin{align*}
(n+3)\;I_{n+2} + (n+1)\;I_n &= \frac{\pi}{2} - \int_0^1 \frac{x^{n+1}(x^2+1)}{1+x^2}\,\mathrm{d}x\\
&= \frac{\pi}{2} - \int_0^1 x^{n+1}\,\mathrm{d}x\\
&= \frac{\pi}{2} - \frac{1}{n+2}
\end{align*}

$$\boxed{(n+3)\;I_{n+2} + (n+1)\;I_n = \frac{\pi}{2} - \frac{1}{n+2}}$$

When $n=2$, $\qquad 5I_4 + 3I_2 = \dfrac{\pi}{2} - \dfrac{1}{4}$

When $n=0$, $\qquad 3I_2 + I_0 = \dfrac{\pi}{2} - \dfrac{1}{2}$

$\implies 3I_2 = \dfrac{\pi}{2} - \dfrac{1}{2} - \left(\dfrac{\pi}{4} - \dfrac{\ln 2}{2}\right) = \dfrac{\pi}{4} - \dfrac{1}{2} + \dfrac{\ln 2}{2}$

$\implies 5I_4 = \dfrac{\pi}{2} - \dfrac{1}{4} - 3I_2 = \dfrac{\pi}{2} - \dfrac{1}{4} - \dfrac{\pi}{4} + \dfrac{1}{2} - \dfrac{\ln 2}{2} = \dfrac{\pi}{4} + \dfrac{1}{4} - \dfrac{\ln 2}{2}$
$$\boxed{I_4 = \frac{\pi}{20} + \frac{1}{20} - \frac{\ln 2}{10}}$$




\bigskip\hrule\bigskip

\noindent
\(
\begin{aligned}[t]
\textbf{(iii)}\ n=1:\ \text{LHS} &= 5I_4\\
\text{RHS} &= A - \frac{1}{2}\;\left( -1 + \frac{1}{2}\right)\\
&= A + \frac{1}{4}
\end{aligned}
\)

$\because 5I_4=\dfrac{\pi}{4} - \dfrac{1}{4} - \dfrac{\ln 2}{2}$\\
$\therefore \text{In order to make LHS=RHS, } A=\dfrac{\pi}{4} - \dfrac{\ln 2}{2}$\\

Assume the equation is true for $n=k$:
$$(4k+1)\;I_{4k}=A-\frac{1}{2}\displaystyle\sum_{r=1}^{2k}\left(-1\right)^{r+1}\frac{1}{2r-1}$$


When $n=k+1$:

Apply (ii) with $n=4k$: \quad $(4k+3)\;I_{4k+2}+(4k+1)\;I_{4k}=\dfrac{\pi}{2}-\dfrac{1}{4k+2}$

\begin{align*}
(4k+3)\;I_{4k+2} &= \frac{\pi}{2}-\frac{1}{4k+2}-(4k+1)\;I_{4k}\\
&= \frac{\pi}{2}-\frac{1}{2(2k+1)}-A+\frac{1}{2}\sum_{r=1}^{2k}(-1)^r\frac{1}{r}\\
&= \frac{\pi}{4}+\frac{\ln 2}{2}-\frac{1}{2(2k+1)}+\frac{1}{2}\sum_{r=1}^{2k}(-1)^r\frac{1}{r}
\end{align*}

$$\because -\dfrac{1}{2(2k+1)}=\dfrac{1}{2}\cdot\dfrac{(-1)^{2k+1}}{2k+1}$$
$$\boxed{\therefore (4k+3)\;I_{4k+2}=\frac{\pi}{4}+\frac{\ln 2}{2}+\frac{1}{2}\sum_{r=1}^{2k+1}(-1)^r\frac{1}{r}}$$

Apply (ii) with $n=4k+2$: \quad $(4k+5)\;I_{4k+4}+(4k+3)\;I_{4k+2}=\dfrac{\pi}{2}-\dfrac{1}{4k+4}$
\begin{align*}
(4k+5)\;I_{4(k+1)} &= \frac{\pi}{2}-\frac{1}{2(2k+2)}-\frac{\pi}{4}-\frac{\ln 2}{2}-\frac{1}{2}\sum_{r=1}^{2k+1}(-1)^r\frac{1}{r}\\
&= \frac{\pi}{4}-\frac{\ln 2}{2}-\frac{1}{2(2k+2)}-\frac{1}{2}\sum_{r=1}^{2k+1}(-1)^r\frac{1}{r}
\end{align*}

$\because (-1)^{2k+2}=1\qquad \therefore-\dfrac{1}{2(2k+2)}=-\dfrac{1}{2}\cdot\dfrac{(-1)^{2k+2}}{2k+2}$
$$(4k+5)\;I_{4(k+1)}=A-\frac{1}{2}\sum_{r=1}^{2(k+1)}(-1)^r\frac{1}{r}$$

By mathematical induction the result is true for all $n\geqslant 1$, with $\boxed{A=\dfrac{\pi}{4}-\dfrac{\ln 2}{2}}$
\bigskip
\hrule

\begin{center}
    This solution was submitted by \textbf{Jon}.\\[4pt]
    Published and refined by \textbf{Jon}.\\[8pt]
    \textit{Last updated: April 7, 2026}
\end{center}

\end{document}

